% =====================================================================
\chapter{Problemas}
\label{cap:problemas}
% =====================================================================

Seis problemas de 10 puntos cada uno (60 puntos en total). La puntuación de cada apartado figura entre paréntesis. Los problemas marcados con \emph{Fuente} están adaptados de tareas, hojas de problemas o libros de texto de cursos universitarios; el enunciado es de redacción propia y los datos se han cambiado para que puedan resolverse con calculadora. Las soluciones completas están en el \cref{ap:soluciones}.

% ---------------------------------------------------------------------
\section*{Problema B.1. El óptimo de Chinchilla a mano (M16)}
\phantomsection\addcontentsline{toc}{section}{Problema B.1. El óptimo de Chinchilla a mano}

\begin{ejercicio}{3}{eje:chinchilla}
Se usa la ley $L(N,D)=E+A/N^{\alpha}+B/D^{\beta}$ con los valores del \cref{tab:datos} y la restricción $C=6ND$, con $C$ fijo.
\begin{enumerate}[label=(\alph*)]
  \item \pt{3} Demuestre que en el óptimo se cumple $\alpha A/N^{\alpha}=\beta B/D^{\beta}$ e interprete la igualdad en una frase.
  \item \pt{2} Deduzca que $\Nopt=G\,(C/6)^{a}$ con $a=\beta/(\alpha+\beta)$ y $G=\bigl(\alpha A/(\beta B)\bigr)^{1/(\alpha+\beta)}$. Calcule $a$ y $G$.
  \item \pt{2} Prediga $\Nopt$ y $\Dopt$ para $C=10^{23}$ y para $C=10^{24}$~FLOP. Para $C=10^{24}$, calcule además $\Dopt/\Nopt$ y la pérdida prevista.
  \item \pt{2} Con $C=10^{24}$, compare con el reparto de la regla de 20 tokens por parámetro ($N$, $D$ y pérdida). ¿Cuál de los dos repartos da menor pérdida según la ley? ¿Por cuánto?
  \item \pt{1} Solo se dispone de $2\cdot 10^{12}$ tokens únicos y no se quiere repetir épocas. Con $C=10^{24}$, ¿qué $N$ se usa y cuánta pérdida se pierde frente al óptimo de~(c)?
\end{enumerate}
\fuente{apartado~(c) adaptado de Stanford CS336, primavera de 2026, tarea~3 («Scaling Laws»), problema \texttt{chinchilla\_isoflops}, apartados (a) y (b): predicción de $\Nopt$ y $\Dopt$ a $10^{23}$ y $10^{24}$~FLOP~\cite{cs336a3}. Aquí se parte de la ley paramétrica en lugar de los perfiles IsoFLOP.}
\end{ejercicio}

% ---------------------------------------------------------------------
\section*{Problema B.2. Estimación de Fermi de un entrenamiento (M16--M17)}
\phantomsection\addcontentsline{toc}{section}{Problema B.2. Estimación de Fermi de un entrenamiento}

\begin{ejercicio}{2}{eje:fermi}
Según su ficha técnica pública, Llama~3 70B se entrenó con unos $1{,}5\cdot 10^{13}$ tokens en GPU H100~\cite{llama3}. Supóngase $N=7\cdot 10^{10}$, 8192 H100, $\MFU=40\,\%$, 700~W por GPU, $\PUE=1{,}2$ y 2{,}5~\$ por GPU-hora (cifras aproximadas e hipotéticas).
\begin{enumerate}[label=(\alph*)]
  \item \pt{1} Calcule el cómputo de entrenamiento $C$ en FLOP.
  \item \pt{3} Calcule las GPU-horas, los días de reloj, la energía en MWh y el coste en dólares.
  \item \pt{2} Un modelo con la forma de GPT-2 XL y la arquitectura de los apuntes (RMSNorm, MLP SwiGLU, RoPE, \emph{embeddings} de entrada y salida no compartidos, sin sesgos) tiene vocabulario $V=50\,257$, contexto de 1024 tokens, $L=48$ capas, $d=1600$, 25 cabezas y $\dff=4288$. Calcule su número de parámetros entrenables y la memoria para cargarlo en fp32. ¿Cuánto tardaría entrenarlo $4\cdot 10^{5}$ pasos con lotes de 1024 secuencias de 1024 tokens en una sola H100 a 495~TFLOP/s (TF32) con $\MFU=50\,\%$? Use $C\approx 6ND$, que ignora la atención.
  \item \pt{2} La ficha pública de Llama~3 70B declara unos $6{,}4\cdot 10^{6}$ GPU-horas. ¿Qué MFU efectiva implica? Dé dos razones por las que es menor que el 40\,\%.
  \item \pt{2} ¿Cuántos tokens por parámetro usa Llama~3 70B? ¿Qué $N$ y $D$ daría la regla de 20 tokens con el mismo $C$? Comparando el coste de servir $10^{13}$ tokens con cada modelo (unos $2N$~FLOP por token), explique por qué tiene sentido «sobreentrenar».
\end{enumerate}
\fuente{apartado~(c) adaptado de Stanford CS336, primavera de 2026, tarea~1 («Basics»), problemas \texttt{transformer\_accounting}~(a) y \texttt{adamw\_accounting}~(d)~\cite{cs336a1}.}
\end{ejercicio}

% ---------------------------------------------------------------------
\section*{Problema B.3. Plan de paralelismo 3D (M17)}
\phantomsection\addcontentsline{toc}{section}{Problema B.3. Plan de paralelismo 3D}

\begin{ejercicio}{3}{eje:paralelismo}
Se entrena un modelo denso de $\Psi=7\cdot 10^{10}$ parámetros y 80 capas en 2048 H100 (256 nodos de 8 GPU) con paralelismo de tensores $\mathrm{TP}=8$, de etapas $\mathrm{PP}=4$ y de datos $\mathrm{DP}=64$, como en la \cref{fig:3d}. El lote global es de 1024 secuencias de 4096 tokens y el microlote, de una secuencia.

\begin{figure}[h]
\centering
\begin{tikzpicture}[font=\small,
  etapa/.style={draw,rounded corners=2pt,minimum width=24mm,minimum height=11mm,align=center}]
  \node[etapa] (s1) at (0,0) {Etapa 1\\capas 1--20};
  \node[etapa] (s2) at (3.0,0) {Etapa 2\\capas 21--40};
  \node[etapa] (s3) at (6.0,0) {Etapa 3\\capas 41--60};
  \node[etapa] (s4) at (9.0,0) {Etapa 4\\capas 61--80};
  \draw[-{Latex}] (s1) -- (s2);
  \draw[-{Latex}] (s2) -- (s3);
  \draw[-{Latex}] (s3) -- (s4);
  \node[draw,dashed,regla,fit=(s1)(s4),inner sep=6pt,label={[font=\footnotesize]above:réplica de datos $1$ de $64$; cada etapa, $\mathrm{TP}=8$ dentro de un nodo (NVLink)}] (rep) {};
  \node[font=\footnotesize,align=center] at (4.5,-1.5) {all-reduce de gradientes entre las 64 réplicas por InfiniBand};
\end{tikzpicture}
\caption{Plan de paralelismo del \cref{eje:paralelismo}.}\label{fig:3d}
\end{figure}

\begin{enumerate}[label=(\alph*)]
  \item \pt{1} Compruebe el número de GPU. ¿Por qué $\mathrm{TP}=8$ y no $\mathrm{TP}=16$?
  \item \pt{2} Calcule la memoria de estado por GPU sin ZeRO (16 bytes por parámetro, repartidos solo entre las $\mathrm{TP}\cdot\mathrm{PP}$ GPU de una réplica). ¿Cabe en 80~GB? ¿Sería viable el paralelismo de datos puro, sin TP ni PP?
  \item \pt{2} Con ZeRO-1 sobre la dimensión de datos (pesos y gradientes repartidos entre $\mathrm{TP}\cdot\mathrm{PP}$; el estado del optimizador, además, entre las $\mathrm{DP}$ réplicas), ¿cuánta memoria de estado queda por GPU?
  \item \pt{2} Calcule los tokens por paso, el número $m$ de microlotes por réplica y la burbuja del pipeline.
  \item \pt{2} Calcule el tiempo de cómputo de un paso ($\MFU=40\,\%$) y el del all-reduce en anillo de los gradientes bf16 de cada GPU entre las 64 réplicas a 50~GB/s. ¿Es la comunicación un cuello de botella? Si se redujera el lote global manteniendo $\mathrm{DP}=64$, ¿con cuántas secuencias por réplica, como mínimo, quedaría la comunicación por debajo del cómputo?
  \item \pt{1} ¿Cuántos pasos y cuántos días hacen falta para $1{,}5\cdot 10^{13}$ tokens?
\end{enumerate}
\fuente{apartado~(b) adaptado de Austin \emph{et al.}, \emph{How to Scale Your Model}, parte~5 («Training»), problemas~2 y~3(a)~\cite{hsym}; apartado~(e) adaptado de Stanford CS336, primavera de 2026, tarea~2 («Systems»), problema \texttt{data\_parallel\_calcs}~(c)~\cite{cs336a2}.}
\end{ejercicio}

% ---------------------------------------------------------------------
\section*{Problema B.4. Mezcla de expertos (M18)}
\phantomsection\addcontentsline{toc}{section}{Problema B.4. Mezcla de expertos}

\begin{ejercicio}{2}{eje:moe}
«Cernícalo-MoE» es un modelo inventado con $d=5120$ y $L=40$ capas. Cada capa tiene 16 expertos con enrutado top-2; cada experto es un MLP SwiGLU de tres matrices sin sesgos con $\dff=8192$. La atención es GQA con 40 cabezas de consulta, $\nkv=8$ y $\dhd=128$. El vocabulario tiene $V=64\,000$ tokens, con \emph{embeddings} de entrada y salida no compartidos. En (a) y (b) se ignoran las normalizaciones y el enrutador.
\begin{enumerate}[label=(\alph*)]
  \item \pt{2} Calcule los parámetros de un experto, de la atención de una capa y de los \emph{embeddings}. Dé el total, los activos por token y la fracción activa.
  \item \pt{2} Calcule los FLOP por token en la pasada hacia delante y la memoria de pesos en bf16. ¿Cuántas H100 de 80~GB hacen falta, como mínimo, solo para los pesos? ¿Qué ahorra una MoE respecto a un modelo denso con los mismos parámetros totales y qué no ahorra?
  \item \pt{2} Para un token, los logits del enrutador son
  \[
  z=(0{,}4;\ 1{,}9;\ -0{,}3;\ 1{,}2;\ 0{,}0;\ 2{,}6;\ 0{,}8;\ -1{,}1;\ 0{,}5;\ 1{,}0;\ -0{,}6;\ 0{,}3;\ 1{,}7;\ -0{,}2;\ 0{,}1;\ 0{,}9).
  \]
  Elija los expertos top-2 y calcule sus pesos $g_i$, con la softmax restringida a los elegidos (M18, §4.1).
  \item \pt{2} En un lote, 4 expertos reciben cada uno una fracción $f_i=0{,}16$ del tráfico con $P_i=0{,}12$, y los otros 12 reciben $f_i=0{,}03$ con $P_i=13/300$. Compruebe que $f$ y $P$ suman~1 y calcule $\Lbal=N\sum_i f_iP_i$. Compárelo con el valor del reparto perfecto.
  \item \pt{2} Con un lote de 8192 tokens, $k=2$ y factor de capacidad $1{,}25$, ¿cuántas plazas tiene cada experto? ¿Cuántas asignaciones descarta cada experto saturado del apartado~(d)?
\end{enumerate}
\fuente{apartados (a) y (b) adaptados de Austin \emph{et al.}, \emph{How to Scale Your Model}, parte~7 («Inference»), problema~5 (variante MoE de un modelo denso: parámetros totales y activos, caché KV y FLOP)~\cite{hsym}.}
\end{ejercicio}

% ---------------------------------------------------------------------
\section*{Problema B.5. Servir un modelo denso (M19)}
\phantomsection\addcontentsline{toc}{section}{Problema B.5. Servir un modelo denso}

\begin{ejercicio}{3}{eje:servicio}
«Halcón-35» es un modelo denso inventado con $L=56$, $d=7168$, 56 cabezas de consulta y GQA con $\nkv=8$ y $\dhd=128$. Se toma $N=3{,}5\cdot 10^{10}$ parámetros. Se sirve en bf16 en 2 H100 con paralelismo de tensores, reservando 10~GB para activaciones y sistema; el ancho de banda agregado es la suma del de las dos GPU, y el rendimiento de cómputo, también.
\begin{enumerate}[label=(\alph*)]
  \item \pt{2} Compruebe que $12Ld^{2}$ se acerca a $N$. Calcule la memoria de pesos, la memoria libre para la caché KV, la caché por token y la caché de una secuencia de 8192 tokens.
  \item \pt{2} ¿Cuántas secuencias de 8192 tokens caben? ¿Y si el modelo usara atención multicabeza clásica, con 56 cabezas de clave y valor?
  \item \pt{3} Calcule el tiempo por paso de decodificación y los tokens/s (totales y por usuario) en dos casos: (i)~lote~1 con 8192 tokens de contexto; (ii)~lote~32 con 4096 tokens de contexto cada uno. Compruebe que en (ii) el cómputo no limita. Ignorando la caché KV, ¿a partir de qué lote pasaría la decodificación a estar limitada por el cómputo? ¿Cabría ese lote, con 4096 tokens por secuencia?
  \item \pt{1} Calcule el tiempo hasta el primer token de un prompt de 4000 tokens con $\MFU=50\,\%$ en el \emph{prefill}.
  \item \pt{2} Con pesos int4 en grupos de 64 con escala fp16, ¿cuál es la cota de tokens/s del caso~(i)? Si además se usa decodificación especulativa con tasa de aceptación $\alpha=0{,}75$, $\gamma=5$ propuestas por ronda y un borrador que cuesta el 4\,\% de una pasada del modelo, ¿qué aceleración adicional cabe esperar?
\end{enumerate}
\fuente{apartados (a)–(c) adaptados de Austin \emph{et al.}, \emph{How to Scale Your Model}, parte~7 («Inference»), problemas~1 a~3 (caché KV por token, lote máximo y tiempo de lectura de los pesos), y parte~1 («Rooflines»), problema~5 (lote crítico de la H100)~\cite{hsym}.}
\end{ejercicio}

% ---------------------------------------------------------------------
\section*{Problema B.6. Extrapolar, deduplicar y desconfiar (M16)}
\phantomsection\addcontentsline{toc}{section}{Problema B.6. Extrapolar, deduplicar y desconfiar}

\begin{ejercicio}{2}{eje:extrapolar}
\begin{enumerate}[label=(\alph*)]
  \item \pt{4} Dos entrenamientos pequeños dan pérdida $3{,}2$ con $C=10^{20}$~FLOP y $2{,}6$ con $C=10^{22}$. Supóngase $L(C)=1{,}7+k\,C^{-b}$. Calcule $b$, prediga $L(10^{24})$ y calcule el cómputo necesario para reducir a la mitad el exceso sobre $1{,}7$ respecto al de $10^{22}$.
  \item \pt{3} Se deduplica con MinHash y LSH por bandas con una firma de 128 valores: $b=16$ bandas de $r=8$ filas. Calcule $P(\text{candidato})=1-(1-s^{r})^{b}$ para $s=0{,}5$, $0{,}7$ y $0{,}9$; el umbral aproximado $(1/b)^{1/r}$, y el umbral exacto $s_{1/2}$, en el que $P=1/2$. ¿Es más o menos estricto que la configuración $r=5$, $b=20$ de los apuntes? ¿Qué tipo de error aumenta?
  \item \pt{3} Una respuesta de $k=20$ tokens solo puntúa si acierta todos. Para $p=0{,}9$, $0{,}95$, $0{,}98$ y $0{,}995$ por token, calcule la exactitud exacta $p^{20}$ y la log-probabilidad de la respuesta correcta, $20\ln p$. Compare la forma de las dos curvas y explique la tesis de Schaeffer, Miranda y Koyejo (2023).
\end{enumerate}
\fuente{apartado~(b) adaptado de Leskovec, Rajaraman y Ullman, \emph{Mining of Massive Datasets}, 3.ª ed., ejercicios~3.4.1 y~3.4.2, libro de texto de Stanford CS246~\cite{mmds,cs246}; el procedimiento (firmas MinHash, bandas LSH y verificación de la similitud de Jaccard de los candidatos) es el del problema \texttt{minhash\_deduplication} de la tarea~4 de Stanford CS336, primavera de 2026~\cite{cs336a4}.}
\end{ejercicio}
